% ---------------------------------------------------------------
%  MNRAS version, converted from the REVTeX source by
%  Plotting_Scripts/to_mnras.py.  Do not edit this file by hand:
%  edit the REVTeX source and re-run the converter, or the two will
%  drift apart.
% ---------------------------------------------------------------
\documentclass[fleqn,usenatbib]{mnras}
\usepackage{graphicx}
\usepackage{amsmath,amssymb,bm}
\usepackage{txfonts}

\newcommand{\Esym}{E_{\mathrm{sym}}}
\newcommand{\Lsym}{L_{\mathrm{sym}}}
\newcommand{\DRnp}{\Delta R_{np}}
\newcommand{\Msun}{M_\odot}
\newcommand{\Teff}{T_{\mathrm{eff}}}
\newcommand{\fmc}{\mathrm{fm}^{-3}}

\title[Magnetized envelopes in NS cooling]{Magnetized envelopes and the symmetry energy from neutron star cooling: a silently inactive field and what happens when it is switched on}

\author[Singh \& Nandi]{
Nisha Singh$^{1}$\thanks{E-mail: ns515@snu.edu.in}
and Rana Nandi$^{2}$
\\
$^{1}$Department of Physics, Shiv Nadar Institution of Eminence, Gautam Buddha Nagar, Uttar Pradesh 201314, India\\
$^{2}$Department of Physics, Indian Institute of Technology Delhi, Hauz Khas, New Delhi 110016, India
}

\begin{document}
\label{firstpage}
\maketitle

\begin{abstract}
The relation between a neutron star's interior temperature and the surface
temperature we observe depends on the magnetic field threading its envelope.
We report two things. First, a defect in a widely used thermal evolution code:
the array holding the surface field is never populated in the distributed
build, so selecting the magnetized iron envelope returns the field-free result
identically rather than approximately, and the field appears to have no effect.
We show this is an exact algebraic identity of the envelope formula at zero
field, which is why it is invisible in testing. Second, with the field actually
applied, we ask what magnetized envelopes do to a symmetry energy constraint
derived from cooling. The sixteen observed neutron stars used for such
constraints span $4.6$ decades in surface dipole field, from
$5.8\times10^{8}$~G to $2.5\times10^{13}$~G, and --- this is the part that
matters --- that spread is not a trend but class structure. The youngest
objects are central compact objects with anomalously weak fields; the strongest
fields belong to two X-ray dim isolated neutron stars in the middle of the
cooling diagram; the two genuinely oldest stars are weakly magnetized again.
Because the field is non-monotonic across the diagram, any single global choice
produces an error whose sign differs from one region to another, and no single
choice removes it. Assigning each star its own measured field improves the fit,
from $\chi^2/N = 1.397$ to $1.335$ --- a change of $0.79\sigma$, which we do not
claim is decisive --- whereas adopting a single global $10^{12}$~G field makes
it \emph{worse}, at $1.467$. The ordering is the result: the field matters, but
only if it is right star by star. The symmetry energy slope inferred from cooling is
unchanged at $\Lsym = 50$--$75$~MeV, which we report as a robustness result for
that constraint.
\end{abstract}

\begin{keywords}
neutron star cooling -- magnetized envelopes -- symmetry energy
\end{keywords}


% =====================================================================
\section{Introduction}

A cooling neutron star is observed through its envelope: the outermost few
metres, which carry the temperature gradient between the nearly isothermal
interior and the radiating surface. Everything a thermal evolution calculation
predicts about an observable temperature passes through the relation
$\Teff(T_b)$ that the envelope supplies, where $T_b$ is the temperature at its
base \citep{GPE1983,HernquistApplegate1984,PCY1997}.

That relation depends on the magnetic field. A field suppresses electron
thermal conduction across field lines and enhances it along them, so the
surface temperature becomes non-uniform and the angle-averaged $\Teff$ at fixed
$T_b$ shifts. Potekhin and Yakovlev \citep{PY2001} give a fitting formula for an
iron envelope threaded by a dipole field, and it is the standard treatment.

Constraints on the nuclear symmetry energy from cooling
\citep{PaperI,Sarkar2023} compare predicted surface temperatures with observed
ones across a population of stars. Such comparisons are sensitive to the
envelope at the level of a few hundredths of a dex, which is the size of the
quoted observational uncertainties. Whether the field is included is therefore
not a detail.

This paper makes two points. Section~\ref{sec:bug} reports that in the
distributed build of the thermal evolution code \textsc{NSCool}
\citep{Page2016}, requesting the magnetized envelope has no effect at all, for a
reason that is exact rather than approximate and therefore does not show up as
a small discrepancy. Section~\ref{sec:fields} onwards asks what including the
field actually does, using each star's own measured dipole field rather than a
single assumed value.

% =====================================================================
\section{A magnetized envelope that was not magnetized}
\label{sec:bug}

\subsection{The mechanism}

In \textsc{NSCool}, the envelope is selected by an integer \texttt{IFTEFF}.
The value 4 selects the magnetized iron relation of \citet{PY2001},
implemented as \texttt{fteff\_field\_iron(Tb, bfield)}. The dispatching routine
passes \texttt{bf\_r(imax)}, the surface element of the radial field array.

That array is never assigned in the distributed build. The only code that sets
it lives in a field-evolution include file which is not part of the compiled
source; the user guide states that field evolution is no longer implemented.
\texttt{bf\_r(imax)} therefore reaches the envelope routine as an
uninitialised common-block value, in practice zero.

The consequence is exact. Writing $B_{12}=B/10^{12}$~G and $T_9=T_b/10^9$~K,
the fitting formula of \citet{PY2001} enters through
\begin{equation}
\beta = 0.074\,\frac{\sqrt{B_{12}}}{T_9^{0.45}} ,
\label{eq:beta}
\end{equation}
and returns
\begin{equation}
\Teff = \Teff^{(0)}
\left[\frac{1 + a_1\beta^2 + a_2\beta^3 + 0.007\,a_3\beta^4}
           {1 + a_3\beta^2}\right]^{1/4} ,
\label{eq:ratio}
\end{equation}
where $\Teff^{(0)}$ is the field-free iron result. The coefficients $a_1$ and
$a_2$ depend on $T_9$ alone; $a_3$ carries in addition a factor
$(1-0.4\,x)^{-1/2}$ with $x = 2GM/Rc^2$ the stellar compactness. Their form
does not matter for what follows. At $B=0$ we have $\beta=0$, so every term in
$\beta$ vanishes, the bracket in Eq.~(\ref{eq:ratio}) is unity
\emph{identically} whatever the $a_i$ are, and $\Teff = \Teff^{(0)}$.

So the magnetized envelope does not merely give a slightly wrong answer at zero
field: it reduces exactly to the unmagnetized envelope. A calculation run with
\texttt{IFTEFF}~$=4$ returns bit for bit what \texttt{IFTEFF}~$=3$ with zero
light-element mass returns, and the natural reading of that output is that the
magnetic field makes no difference to cooling.

\subsection{Demonstration}

We verified this directly. Running the same equation of state at
$1.4\,\Msun$ under \texttt{IFTEFF}~$=4$ with $B=0$ and under
\texttt{IFTEFF}~$=3$ with $\eta=0$ gives time grids and surface temperatures
that agree to the last stored digit: the maximum difference in $\Teff$ over the
294 output steps is $0.000$~K, and the two runs differ only by one banner line
that the field-free branch prints and the magnetized branch does not. Assigning
the field array and repeating the $B=10^{12}$~G case then produces a difference,
as it must.

Both statements are enforced as a regression test, so that a rebuild from a
clean upstream tree cannot silently reinstate the defect.

\begin{figure}[t]
\includegraphics[width=\columnwidth]{figs/fig_bug_proof.pdf}
\caption{The defect, shown rather than asserted. Upper panel: cooling tracks for
one equation of state at $1.4\,\Msun$. The magnetized envelope evaluated at
$B=0$ (thin red) lies exactly on the field-free track (thick grey); the two are
indistinguishable because they are identical. With the field array populated and
$B=10^{12}$~G (blue) the track separates. Lower panel: the differences. The
$B=0$ residual is identically zero at every one of the 294 output steps, which
is why the defect produces no small discrepancy for testing to catch.}
\label{fig:bug}
\end{figure}

\subsection{Scope}

We state the scope carefully. This is a property of the distributed build, not
an error in the physics of \citet{PY2001}, whose formula is correctly
transcribed. Nor is it a criticism of a code that has been made freely
available and is widely useful. But any published result that selected the
magnetized iron envelope in this build, and reported that the field had little
effect, should be re-examined: the field was not applied.

% =====================================================================
\section{The observed stars span nearly five decades in field}
\label{sec:fields}

\subsection{Where the fields come from}

For the thirteen pulsating objects in the sample of \citet{PaperI} we take
$P$ and $\dot P$ from version 2.6.0 of the ATNF Pulsar Catalogue
\citep{ATNF2005} and form the characteristic surface dipole field
\begin{equation}
B = 3.2\times10^{19}\sqrt{P\dot P}~\mathrm{G},
\label{eq:bsurf}
\end{equation}
which is the definition the catalogue itself tabulates. We recomputed rather
than copied, and checked the result against published values for four objects
spanning the range.

Three objects have no spin-down field. The neutron star in Cassiopeia~A and
XMMU~J173203.3$-$344518 show no pulsations at all
\citep{HoHeinke2009,Klochkov2013}, so no $\dot P$ exists; RX~J0822$-$4300
pulsates at $P=0.112$~s but its $\dot P$ is not determined well enough to give
a field. All three are central compact objects, a class whose members have
anomalously weak fields --- ``anti-magnetars'', $B \lesssim 10^{11}$~G
\citep{DeLuca2017}. We assign them the field-free track and show below that
nothing sensitive rides on that choice, because the envelope shift below
$10^{11}$~G is small.

\begin{table}[t]
\caption{Surface dipole fields, from Eq.~(\ref{eq:bsurf}) with $P$ and
$\dot P$ from \citet{ATNF2005}. The three central compact objects have no
measured $\dot P$ and are treated as anti-magnetars \citep{DeLuca2017}.
Ages and temperatures are as in \citet{PaperI}.}
\label{tab:fields}

\begin{tabular}{llcc}
\hline
Source & ATNF name & $\log_{10}t$ & $\log_{10}(B/\mathrm{G})$\\
\hline
Cas\,A      & ---          & 2.51 & $\lesssim 11$ \\
XMMU\,J1732 & ---          & 3.48 & $\lesssim 11$ \\
RX\,J0822   & ---          & 3.57 & $\lesssim 11$ \\
1E\,1207    & J1210$-$5226 & 3.85 & 10.99 \\
PSR\,B0833  & J0835$-$4510 & 4.05 & 12.53 \\
PSR\,B1706  & J1709$-$4429 & 4.24 & 12.50 \\
PSR\,J0538  & J0538$+$2817 & 4.48 & 11.87 \\
PSR\,B2334  & J2337$+$6151 & 4.61 & 13.00 \\
PSR\,B0656  & J0659$+$1414 & 5.05 & 12.67 \\
Geminga     & J0633$+$1746 & 5.53 & 12.21 \\
RX\,J1856   & J1856$-$3754 & 5.70 & 13.17 \\
PSR\,B1055  & J1057$-$5226 & 5.73 & 12.03 \\
RX\,J0720   & J0720$-$3125 & 5.78 & 13.39 \\
PSR\,J2043  & J2043$+$2740 & 5.88 & 11.55 \\
PSR\,J0437  & J0437$-$4715 & 6.85 & 8.76  \\
PSR\,B0950  & J0953$+$0755 & 7.24 & 11.39 \\
\hline
\end{tabular}

\end{table}

\subsection{Why the spread matters, and why it is not a trend}

Table~\ref{tab:fields} spans $4.6$ decades, from $5.8\times10^{8}$~G for the
recycled pulsar PSR~J0437$-$4715 to $2.5\times10^{13}$~G for the X-ray dim
isolated neutron star RX~J0720$-$3125. Eight of the thirteen measured fields
exceed $10^{12}$~G.

It would be convenient if the field simply grew or decayed with age, because
then its effect could be absorbed into a redefinition of the age axis. It does
not. Over the thirteen measured objects the rank correlation between
$\log B$ and $\log t$ is $\rho = -0.19$ with $p = 0.54$: there is no monotonic
trend to speak of. What there is instead is \emph{class structure}, and the
classes sit in different places on the cooling diagram:

\begin{center}
\begin{tabular}{lccl}
\hline
class & $n$ & $\log_{10}t$ & $B$ (G)\\
\hline
central compact objects & 3 & 2.5--3.6 & $\lesssim 10^{11}$ (assumed)\\
ordinary pulsars        & 9 & 3.9--5.9 & $9.8\times10^{10}$--$9.9\times10^{12}$\\
X-ray dim isolated NS   & 2 & 5.7--5.8 & $1.5$--$2.5\times10^{13}$\\
the two oldest          & 2 & 6.9--7.2 & $5.8\times10^{8}$--$2.4\times10^{11}$\\
\hline
\end{tabular}
\end{center}

Read down that table and the field goes weak, then strong, then strongest,
then weak again. The two X-ray dim objects carry the strongest fields in the
sample but are neither the oldest nor the coolest members of it; the genuinely
oldest two, PSR~J0437$-$4715 and PSR~B0950$+$08, carry the weakest measured
fields.

This non-monotonicity is what makes a single global field inadequate, and the
argument is stronger than a simple trend would have given. A monotonic trend
could be compensated, at least partly, by shifting one of the axes. A field
that is weak among the young objects, strongest in the middle of the diagram
and weak again at the old end cannot be: whatever global value is chosen, the
resulting error has \emph{different signs} in different parts of the very
diagram being fitted. No single choice removes it.

That is the argument for giving each star its own field rather than adopting
one value for all of them.

\begin{figure}[t]
\includegraphics[width=\columnwidth]{figs/fig_field_landscape.pdf}
\caption{Surface dipole field against age for the sixteen stars, by class. The
three central compact objects have no measured $\dot P$ and are shown as upper
limits. The field is weak among the youngest objects, reaches its maximum in the
two X-ray dim isolated neutron stars in the middle of the diagram, and is weak
again in the two genuinely oldest stars. There is no monotonic trend
($\rho=-0.19$, $p=0.54$), which is why no single global field can absorb the
variation: the residual would change sign across the diagram.}
\label{fig:landscape}
\end{figure}

% =====================================================================
\section{Method}

Everything except the envelope is held identical to \citet{PaperI}: the
same 864 equations of state, the same five parametrizations with only the
isovector couplings varied, vortex creep heating so that all sixteen stars are
reachable without extrapolation, carbon envelopes for the two objects with
published carbon atmospheres, the corrected RX~J0822 temperature, and a
$1.4\pm0.3\,\Msun$ mass prior. The single change is that the field-free iron
envelope is replaced by the magnetized one of \citet{PY2001}.

Cooling curves are computed on a grid of
$B = 0,\,10^{11},\,10^{12},\,10^{13},\,10^{14}$~G at each of
$M = 1.0,\,1.4,\,1.8,\,2.0\,\Msun$. For a star of field $B$, the predicted
$\log_{10}\Teff$ at its age is interpolated linearly in $\log_{10}B$ between
the two bracketing grid values; stars with $B \le 10^{11}$~G take the
field-free track directly.

We report three fits, so that the effect is attributable rather than merely
present:
\begin{enumerate}
\item $B=0$ for every star, which reproduces \citet{PaperI} and serves as
the check that the pipeline is unchanged;
\item $B=10^{12}$~G for every star, the naive global choice;
\item each star's own measured field, Table~\ref{tab:fields}.
\end{enumerate}

% =====================================================================
\section{Results}

\subsection{The size of the envelope effect}

Table~\ref{tab:shift} gives the shift the field produces in the observable
surface temperature, at fixed age, relative to the field-free envelope.

\begin{table}[t]
\caption{Change in $\log_{10}\Teff$ at fixed age when the field-free iron
envelope is replaced by the magnetized one, at $1.4\,\Msun$. Each entry is the
mean over \emph{all 864} equations of state, not a single representative curve;
the spread across the 864 is at most $0.02$~dex and is largest at
$10^{14}$~G. The effect is \emph{not} monotonic in $B$.}
\label{tab:shift}

\begin{tabular}{ccccc}
\hline
 & \multicolumn{4}{c}{$\Delta\log_{10}\Teff$ (dex)}\\
$\log_{10}t$ & $10^{11}$~G & $10^{12}$~G & $10^{13}$~G & $10^{14}$~G\\
\hline
2 & $-0.014$ & $-0.030$ & $-0.019$ & $+0.022$\\
3 & $-0.020$ & $-0.032$ & $-0.014$ & $+0.036$\\
4 & $-0.022$ & $-0.031$ & $-0.010$ & $+0.043$\\
5 & $-0.020$ & $-0.025$ & $-0.004$ & $+0.040$\\
6 & $+0.014$ & $+0.016$ & $-0.001$ & $-0.030$\\
7 & $+0.001$ & $+0.000$ & $-0.004$ & $-0.013$\\
\hline
\end{tabular}

\end{table}

Two features deserve comment. First, the magnitude: through the
neutrino-cooling era the shift reaches $-0.032$~dex at $10^{12}$~G and
$+0.043$~dex at $10^{14}$~G. That is a third to nearly a half of the
observational uncertainty on a typical object in the sample ($\sigma = 0.10$~dex),
and comparable to the uncertainty on the best-measured ones ($0.05$~dex). It is
not negligible at the precision these fits claim.

Second, and less obviously, \textbf{the shift is not monotonic in $B$}. It
deepens from $10^{11}$ to $10^{12}$~G, becomes shallower by $10^{13}$~G, and has
changed sign by $10^{14}$~G. This holds for every one of the 864 equations of
state, not merely on average. This follows from the anisotropy of the
conduction. Potekhin and Yakovlev \citep{PY2001} write the effective
conductivity of the blanketing envelope as
$\kappa = \kappa_\parallel\cos^2\theta + \kappa_\perp\sin^2\theta$, with
$\theta$ the angle between the field lines and the surface normal, so the heat
flow is carried by $\kappa_\parallel$ where the field is radial (near the
magnetic poles) and by $\kappa_\perp$ where it is tangential (near the
equator). A radial field alone ``would always lower the thermal insulation of
the blanketing envelope'' \citep{PY2001}; for a dipole the polar and equatorial
regions pull in opposite directions and partly cancel, which is why the same
authors record that the field ``did not necessarily accelerate the cooling''.
Our angle-averaged result turns over because that cancellation is not
field-independent. The practical consequence is that one
cannot capture the effect by applying a single correction factor, or by taking a
representative field for the population --- which is what the next section
measures.

\begin{figure}[t]
\includegraphics[width=\columnwidth]{figs/fig_shift_vs_B.pdf}
\caption{The envelope shift against field strength at four ages, at
$1.4\,\Msun$. Points are the mean over all 864 equations of state and the bands
give the full range across them. The shift deepens to $10^{12}$~G, becomes
shallower by $10^{13}$~G and reverses sign by $10^{14}$~G. Every one of the 864
equations of state follows this pattern; it is not an average of differing
behaviours.}
\label{fig:shift}
\end{figure}

\subsection{Effect on the fit and on the symmetry energy}

\begin{table}[t]
\caption{Fit quality and the inferred symmetry energy slope under three
treatments of the envelope field. All three use the same 864 equations of
state, the same sixteen stars, and differ only in the field assigned to each
star. The entire field grid is computed: no interpolation in this table falls
back on a neighbouring grid value.}
\label{tab:fits}

\begin{tabular}{lccc}
\hline
Field treatment & $\chi^2/N$ & $\chi^2/\mathrm{dof}$ & $\Lsym$ (68\%)\\
\hline
$B=0$ for every star \citep{PaperI}    & 1.397 & 1.596 & $[50,75]$\\
$B=10^{12}$~G for every star          & 1.467 & 1.677 & $[50,75]$\\
each star's measured $B$ (this work)  & 1.335 & 1.525 & $[50,75]$\\
\hline
\end{tabular}

\end{table}

Table~\ref{tab:fits} is the central result, and the ordering in it is the point
rather than any single entry.

\textbf{A single global field makes the fit worse.} Adopting $10^{12}$~G --- a
reasonable-looking choice, close to the sample median of $1.6\times10^{12}$~G ---
degrades $\chi^2/N$ from $1.397$ to $1.467$. Including the field as a uniform
property of the population is therefore not merely imprecise; it is worse than
omitting it.

\textbf{The measured fields make it better.} Giving each star the field
inferred from its own spin-down improves $\chi^2/N$ to $1.335$, a change of
$-0.99$ in total $\chi^2$, or $0.79\sigma$ after inflating uncertainties by the
Particle Data Group scale factor \citep{PDG2024}.

The two statements together are the substance of this paper. The envelope field
is a real effect of the size shown in Table~\ref{tab:shift}, but it enters the
population fit \emph{only} through the star-to-star variation, and that
variation is structured by object class rather than random. Averaging it away
discards the information and adds a systematic in its place.

We stress the modest size of the improvement. At $0.79\sigma$ this is not a
decisive preference, and we do not claim that including fields is demanded by
the data. What the data do show is a consistent ordering: measured fields
better than no field, and no field better than one global field.

\subsection{The symmetry energy constraint is unaffected}

\textbf{The $\Lsym$ interval is $[50,75]$~MeV in all three treatments.} The
constraint reported in \citet{PaperI} therefore survives the inclusion of
magnetized envelopes with per-star fields, unchanged to the resolution of the
grid.

This is worth stating positively rather than as an absence. The envelope field
is exactly the kind of systematic that could have masqueraded as a symmetry
energy effect, because it varies systematically across the cooling diagram and
because the symmetry energy acts by tilting the cooling track against age. It
does not: the shift it produces is absorbed almost entirely into the fit quality
rather than into the inferred parameter. The cooling constraint on $\Lsym$ is
robust to it.

% =====================================================================
\section{Systematic uncertainties}

\paragraph{Dipole geometry.} The relation of \citet{PY2001} is for a
dipole field and returns an angle-averaged $\Teff$. Real fields have
higher multipoles, and the observed temperature depends on viewing geometry.

\paragraph{Field evolution.} We hold $B$ constant over the star's life. Fields
are expected to decay, so the value inferred from present-day spin-down is a
lower bound on the field at early times for a decaying field, and the shift we
compute is correspondingly conservative at young ages.

\paragraph{The characteristic field is not the field.} Equation~(\ref{eq:bsurf})
assumes magnetic dipole braking with a fixed moment of inertia and orthogonal
rotator geometry. It is standard, and it is what the literature quotes, but it
is an estimate.

\paragraph{The anti-magnetar assumption.} Three objects are assigned
$B \lesssim 10^{11}$~G on class grounds rather than measurement. Replacing that
assumption with an ordinary $10^{12}$~G field for all three raises the best
$\chi^2$ from $21.35$ to $23.42$, i.e.\ $\chi^2/N$ from $1.335$ to $1.464$, a
change of $1.16\sigma$. The adopted assumption is thus the one the data prefer,
which is a consistency check rather than a fit: the anti-magnetar
classification comes from the literature \citep{DeLuca2017}, not from this
$\chi^2$. Note also that two of the three keep a field-free carbon envelope in
either case, so this test moves RX~J0822 alone.

\paragraph{Envelope composition and field are not independent.} The two carbon
atmosphere objects are treated with a field-free carbon envelope, since a
magnetized carbon relation is not available in this build. Both are
anti-magnetars, so this is consistent, but it is an assumption.

% =====================================================================
\section{Conclusions}

We report three things.

First, in the distributed build of \textsc{NSCool} the array carrying the
surface magnetic field is never populated, and the envelope formula of
\citet{PY2001} reduces \emph{exactly} to the field-free case at zero field.
Selecting the magnetized iron envelope therefore returns the unmagnetized
answer bit for bit. The failure is silent and produces no small discrepancy that
testing would catch; it simply looks like a magnetic field that does not matter.
Any result obtained with that build which reports the field to be unimportant
should be re-examined, because the field was not applied.

Second, the observed cooling population spans $4.6$ decades in surface dipole
field, and that spread is organised by object class rather than by a trend with
age: the youngest objects are anti-magnetar central compact objects, the
strongest fields belong to two X-ray dim isolated neutron stars lying mid-way
along the cooling diagram, and the two genuinely oldest stars are weakly
magnetized. A single global field is consequently not an adequate treatment,
and because the variation is non-monotonic it cannot be absorbed by rescaling
an axis either. We find it is worse than none: adopting $10^{12}$~G for
every star degrades the fit relative to omitting the field entirely, while
using each star's measured field improves it, modestly, by $0.79\sigma$. It is
the ordering of the three, not the size of any one change, that carries the
conclusion.

Third, and reassuringly, the symmetry energy slope inferred from thermal
evolution is unaffected. $\Lsym = 50$--$75$~MeV at 68\% credibility under every
field treatment tested. Since the envelope field varies systematically across
the cooling diagram, and since the symmetry energy acts on the cooling track
through the same variable, this robustness is not automatic and is worth
establishing explicitly.


\section*{Acknowledgements}

N.S. thanks Kenji Nishiwaki for supervision and support, and acknowledges Shiv
Nadar Institution of Eminence for a research fellowship. We thank Dany Page for
making \textsc{NSCool} publicly available. This work made use of the ATNF
Pulsar Catalogue.

\section*{Data availability}

The equation-of-state tables, cooling outputs and analysis products that
underlie this work are available from the authors on reasonable request.
The thermal evolution code is publicly available from its author and is
not redistributed here.

\bibliographystyle{mnras}
\bibliography{refs}

\label{lastpage}
\end{document}
